\subsection{问题一分析}
2025年的1.8倍总量来自观测。若把四大家鱼合并为单一消费者，并把水草、浮游植物、浮游动物和底栖饵料压成一个资源总量，这一总量仍能匹配，但无法区分"鱼类资源量增加"背后的"鱼类增长的驱动机制"和"基础资源的承载压力变化"，也难以解释"水下荒漠"为何同时出现。因此，食性配对必须留在模型中，每类鱼的结果解释分别落到对应资源轨迹及其承压层级。

题目给出的两类观测数据需差异化使用：2025年四大家鱼1.8倍恢复量对应全流域平均口径，可用于约束系统总量参数；湖北江段单网次渔获量提升120%具有明显局地性和网具特异性，更适合作为外部独立验证，不宜与全流域总量数据共同参与参数标定。据此，我们先在统一总量尺度下识别禁渔的直接生态效应，再将局地捕获差异纳入结果分析。

